\begin{frame}{실전적 상한 + ``평가만 하면 안 된다''는 반례}
  단조 증가는 ``나빠지지 않는다''지만 ``얼마나 빨리 좋아지는가''
  는 주지 않음. 수축 성질(4.3절)로부터:
  \[
  \max_s \big|V^{\pi'}(s) - V^\pi(s)\big|
  \le \frac{1 - \gamma}{\gamma}\, \max_s \big|V^{*}(s) -
  V^\pi(s)\big|
  \]
  \begin{block}{실무적 메시지}
    \textbf{초기 정책이 아무리 엉망이어도, $\gamma$ 가 1에서
    충분히 멀지만 한 한 바퀴(평가+개선)에 오차가 \textbf{고정
    비율}($\gamma$ 배)로 줄어든다.} $\gamma = 0.99$ 처럼 1에
    가까우면 상한도 느려져 전체 비용이 크게 불어남.
  \end{block}
  \vskip0.4em
  \begin{block}{``평가만 하면 안 된다''는 반례}
    ``왼쪽'' 정책($V = -10$)을 \textbf{평가만} 하면 ``가치
    $-10$ 이다''가 \textbf{전부}. 개선 단계를 \textbf{한 번만}
    끼워도 \textbf{칸 3이 right로} 바뀌고($V(3): -10 \to -1$),
    그 정보가 다음 평가에서 \textbf{칸 2로, 다음에 칸 1로}
    퍼짐. \textbf{평가 alone}은 점수만 찍고, \textbf{개선
    단계}가 점수표를 ``갈아탈 방향''으로 번역 --- 두 단계의
    \textbf{곱}(점수 측정 × 방향 추출)이 정책 반복의 힘.
  \end{block}
\end{frame}
